% !Mode:: "TeX:UTF-8"
%
% 第二章：预备知识
%
% 本节展示的用法：
%   - 算法环境（algorithm2e）
%   - 表格（table + tabular + booktabs）
%   - 列表（itemize）
%   - 简单公式
%

\chapter{预备知识}
\echapter{Preliminaries}
\label{chap:preliminaries}

本章介绍后续所需的基础知识，包括优化方法和矩阵运算等。

\section{优化方法}
\esection{Optimization Methods}
\label{sec:optimization}

\subsection{梯度下降法}
\label{subsec:gd}

考虑无约束优化问题
\begin{equation}
\min_{x \in \mathbb{R}^n} f(x),
\label{eq:opt_prob}
\end{equation}
梯度下降的迭代格式为 \(x_{k+1} = x_k - \alpha_k \nabla f(x_k)\)。

常用的优化方法包括：
\begin{itemize}
\item \textbf{梯度下降}：一阶迭代方法，收敛慢但每步计算量小。
\item \textbf{牛顿法}：利用二阶 Hessian 矩阵加速收敛。
\item \textbf{拟牛顿法}：近似 Hessian 矩阵，平衡收敛速度与计算开销。
\end{itemize}

算法\ref{alg:sgd}给出了随机梯度下降的实现流程。

\begin{algorithm}[h]
\caption{随机梯度下降（SGD）}
\label{alg:sgd}
\SetAlgoLined
\KwIn{初始参数 \(x_0\)，学习率 \(\{\alpha_k\}\)，迭代次数 \(K\)}
\KwOut{优化结果 \(x_K\)}
\For{\(k = 0\) \KwTo \(K-1\)}{
    采样小批量 \(\mathcal{B}_k\)\;
    计算梯度 \(g_k = \frac{1}{|\mathcal{B}_k|} \sum_{i \in \mathcal{B}_k} \nabla f_i(x_k)\)\;
    更新 \(x_{k+1} = x_k - \alpha_k g_k\)\;
}
\Return \(x_K\)\;
\end{algorithm}

\section{矩阵运算}
\esection{Matrix Operations}
\label{sec:matrix}

表\ref{tab:methods}对比了不同优化方法的特点。

\begin{table}[h]
\centering
\caption{不同优化方法的对比}
\label{tab:methods}
\begin{tabular}{lccc}
\toprule
方法 & 收敛速度 & 内存开销 & 实现难度 \\
\midrule
梯度下降 & 慢 & 低 & 简单 \\
牛顿法 & 快 & 高 & 中等 \\
拟牛顿法 & 较快 & 中等 & 较复杂 \\
\bottomrule
\end{tabular}
\end{table}
